\documentclass{article}
\usepackage[preprint]{neurips_2026}
\newif\ifsurveychinese
\surveychinesefalse
\usepackage{iftex}
\ifPDFTeX
  \usepackage[T1]{fontenc}
  \usepackage[utf8]{inputenc}
\else
  \usepackage{fontspec}
  \setmainfont{TeX Gyre Termes}
\fi
\ifsurveychinese
  \usepackage[UTF8,fontset=fandol,scheme=plain]{ctex}
\fi
\usepackage{amsmath,amssymb,booktabs,tabularx,graphicx}
\usepackage{microtype,url,hyperref}
\title{Survey Title}
\author{}
\begin{document}
\maketitle
\begin{abstract}
% Context, scope, evidence base, synthesis and limitations.
\end{abstract}
\section{Introduction}\label{sec:introduction}
% Motivation, previous surveys, scope, contributions and reading guide.
% Include documented search/screening methods in a short subsection.
% Choose ONE title: Background OR Preliminaries. Do not print both.
\section{Background}\label{sec:background}
% Define necessary tasks, notation and foundations with citations.
% Choose ONE title: Taxonomy OR Unified Framework. Do not print both.
\section{Taxonomy}\label{sec:taxonomy}
% Explain classification criteria, boundaries and links to topic chapters.
% Expand this topic block once per topic in the initial reconnaissance-based
% Survey Outline recorded in scope. There is no fixed topic count.
\section{Topic Name}\label{sec:topic-slug}
% Synthesize problems, methods, progression, conditions and limitations.
\section{Comparative Analysis and Evaluation}\label{sec:comparison}
% Explicit task/data/metric/budget conditions, coverage, conflicts and limitations.
\section{Challenges and Future Directions}\label{sec:challenges}
% Link challenges to evidenced gaps; distinguish inference from source claims.
\section{Conclusion}\label{sec:conclusion}
% Supported insights and boundaries, without new evidence.
\bibliographystyle{plainnat}
\bibliography{references}
\end{document}
